% !TEX root = ../Main.tex
\chapter{内切球}

\textbf{内切球}是平面里\textbf{内切圆}的三维版本：球面内切于多面体的每一个面。它的刻画从内切圆平移而来——球心到每个面的距离都相等，这个公共距离就是半径。

\section{内切圆与内切球}

\state{两个平行的刻画}{
    \begin{itemize}
        \item \textbf{内切圆}：圆心到各边的距离都相等——关心的是\textbf{点线距}。
        \item \textbf{内切球}：球心到各面的距离都相等——关心的是\textbf{点面距}。
    \end{itemize}
}

\begin{center}
\begin{tikzpicture}[scale=1]
    \coordinate (P) at (0.2,2.3);
    \coordinate (A) at (-1.7,-0.2);
    \coordinate (B) at (0.9,-0.6);
    \coordinate (C) at (1.6,0.2);
    \draw[thick] (P) -- (A) -- (B) -- cycle;
    \draw[thick] (P) -- (B);
    \draw[thick, dashed] (A) -- (C) -- (B);
    \draw[thick, dashed] (P) -- (C);
    \coordinate (O) at (0.1,0.3);
    \draw[blue] (O) circle (0.5);
    \fill (O) circle (1pt) node[right] {\small $O$};
    \draw[dashed] (O) -- (0.05,-0.36) node[midway,left]{\small $r$};
    \node at (P) [above] {$P$};
    \node at (A) [below left] {$A$};
    \node at (B) [below] {$B$};
    \node at (C) [right] {$C$};
    \node at (0,-1.4) {\small 内切球：$O$ 到各面距离 $=r$};
\end{tikzpicture}
\qquad
\begin{tikzpicture}[scale=1]
    \coordinate (A) at (-1,2);
    \coordinate (B) at (-1.9,-0.6);
    \coordinate (C) at (1.9,0.1);
    \draw[thick] (A) -- (B) -- (C) -- cycle;
    \coordinate (I) at (-0.2,0.3);
    \draw[blue] (I) circle (0.62);
    \fill (I) circle (1pt) node[above right] {\small $O$};
    \draw[dashed, red] (I) -- ($(B)!(I)!(C)$) node[midway,below right]{\small $r$};
    \draw[dashed, red] (I) -- ($(A)!(I)!(B)$);
    \draw[dashed, red] (I) -- ($(A)!(I)!(C)$);
    \node at (A) [above] {$A$};
    \node at (B) [below left] {$B$};
    \node at (C) [right] {$C$};
    \node at (0,-1.4) {\small 内切圆：$O$ 到各边距离 $=r$};
\end{tikzpicture}
\end{center}

两者都靠\textbf{垂直关系}——半径沿切点垂直于边（或面）。这个垂直关系在求面积或体积时很有用，关键的招法就是\textbf{分割}。

\section{分割法}

\state{内切圆：由面积反推 $r$}{
    把 $\triangle ABC$ 沿内心 $O$ 切成三个小三角形 $\triangle OBC,\triangle OCA,\triangle OAB$，每个小三角形以一条边为底、高都是 $r$。记周长 $C=a+b+c$：
    \[
    S=\frac12 cr+\frac12 ar+\frac12 br=\frac12(a+b+c)r=\frac12 Cr.
    \]
}

\state{内切球：由体积反推 $r$}{
    把多面体沿球心 $O$ 切成若干个"以一个面为底、$O$ 为顶"的小棱锥，每个小棱锥的高都是 $r$。记表面积 $S_{\text{表}}$：
    \[
    V=\sum \frac13 S_i\,r=\frac13\Big(\sum S_i\Big)r=\frac13 S_{\text{表}}\,r.
    \]
    把二维的 $S=\tfrac12 Cr$ 升一维，就是三维的 $V=\tfrac13 S_{\text{表}}r$。中学里最常用它来\textbf{求内切球半径} $r=\dfrac{3V}{S_{\text{表}}}$。
}

\section{哪些立体有内切球？}

\rmkb{
    \begin{enumerate}
        \item 所有\textbf{三棱锥}都有内切球。\checkmark
        \item 并非所有\textbf{四棱锥}都有内切球。$\times$
        \item 一般\textbf{棱柱}不一定有内切球（如底面 $ABCD$ 为矩形而非正方形的直棱柱）。$\times$
    \end{enumerate}
}

\section{例题}

\ex{
    已知四棱锥 $P$-$ABCD$，平面 $PAB\perp$ 平面 $ABCD$，底面 $ABCD$ 为矩形，$\triangle PAB$ 是边长为 $2$ 的正三角形。若 $P$-$ABCD$ 存在内切球，则此内切球的半径为 \underline{\hspace{3em}}，此四棱锥的体积为 \underline{\hspace{3em}}。
}
